% =====================================================================
% 03-Holonomy.tex
% Geometric Hallucination Taxonomy — Mathematical Spine
% Borealis S. Hedling — 2026
% =====================================================================

\section*{03. Holonomy Operator}

This module formalizes parallel transport, holonomy, and drift-based
instability criteria used to define Holonomy Drift Hallucinations (HDH)
in the Geometric Hallucination Taxonomy (GHT). All definitions rely on
the geometric foundations established in 00-Foundations.tex.

% ---------------------------------------------------------------------
\subsection*{3.1 Parallel Transport Along Curves}

Let $\gamma : [0,1] \to M$ be a smooth curve. A vector field
$V(t) \in T_{\gamma(t)}M$ is said to be parallel along $\gamma$ if
\[
\nabla_{\dot{\gamma}(t)} V(t) = 0.
\]

Given an initial vector $V(0) \in T_{\gamma(0)}M$, the parallel transport
equation has a unique solution $V(t)$ for all $t \in [0,1]$.

Parallel transport defines a linear map
\[
P_\gamma : T_{\gamma(0)}M \to T_{\gamma(1)}M,
\]
called the parallel transport operator along $\gamma$.

% ---------------------------------------------------------------------
\subsection*{3.2 Closed Loops and Holonomy}

If $\gamma$ is a closed loop, i.e.
\[
\gamma(0) = \gamma(1),
\]
then parallel transport defines a linear endomorphism
\[
P_\gamma : T_{\gamma(0)}M \to T_{\gamma(0)}M.
\]

The holonomy operator is defined as
\[
H_\gamma = P_\gamma - I,
\]
where $I$ is the identity map on $T_{\gamma(0)}M$.

Holonomy measures the failure of parallel transport to return a vector
to its original direction after traversing a closed loop.

% ---------------------------------------------------------------------
\subsection*{3.3 Holonomy Group}

The holonomy group at $p \in M$ is
\[
\mathrm{Hol}(p)
= \{ P_\gamma : \gamma \text{ is a loop based at } p \}.
\]

This group is a Lie subgroup of $\mathrm{GL}(T_p M)$ and encodes the
global geometric structure of the manifold.

In the context of GHT, holonomy drift corresponds to instability in
parallel transport around latent loops.

% ---------------------------------------------------------------------
\subsection*{3.4 Holonomy Norm}

Given a holonomy operator $H_\gamma$, we define its norm by
\[
\|H_\gamma\| = \sup_{\|v\|_g = 1} \|H_\gamma v\|_g.
\]

This norm measures the magnitude of rotational or shearing drift induced
by the loop $\gamma$.

% ---------------------------------------------------------------------
\subsection*{3.5 Holonomy Drift Instability}

A holonomy drift instability occurs when
\[
\|H_\gamma\| > h_{\mathrm{thr}}.
\]

Directional asymmetry occurs when
\[
\|H_\gamma\| \neq \|H_{\gamma^{-1}}\|,
\]
where $\gamma^{-1}$ is the loop traversed in reverse.

Oscillatory drift occurs when successive loops satisfy
\[
H_{\gamma_{t+1}} \cdot H_{\gamma_t} < 0
\]
in an appropriate operator sense (e.g., alternating dominant eigenvalues).

These behaviors correspond to distinct HDH subtypes.

% ---------------------------------------------------------------------
\subsection*{3.6 Holonomy Along Latent Trajectories}

Let $x(t)$ be a latent trajectory. A closed loop may be formed by
concatenating segments of the trajectory:
\[
\gamma = x|_{[t_0,t_1]} \ast x|_{[t_1,t_2]} \ast \cdots \ast x|_{[t_k,t_0]}.
\]

Holonomy drift along such loops indicates instability in the latent
reasoning path.

In particular:
\begin{itemize}
\item rotational drift corresponds to semantic rotation,
\item shearing drift corresponds to representational distortion,
\item asymmetric drift corresponds to directional instability.
\end{itemize}

These geometric interpretations motivate the HDH classification.

% ---------------------------------------------------------------------
\subsection*{3.7 Holonomy Stability Region}

Define the holonomy stability region
\[
\Omega_h
= \{ p \in M : \|H_\gamma\| \le h_{\mathrm{thr}}
   \text{ for all loops } \gamma \text{ based at } p \}.
\]

A latent trajectory is holonomy–stable if $x(t) \in \Omega_h$ for all $t$.

% ---------------------------------------------------------------------
\subsection*{3.8 Discrete Approximation}

In numerical settings, parallel transport is approximated using discrete
updates:
\[
V_{t+1}
= V_t - \Gamma(x_t)(\dot{x}_t, V_t),
\]
where $\Gamma$ denotes Christoffel symbols.

Holonomy is computed by composing discrete transport operators around
synthetic loops.

The Python toolkit implements this approximation and validates it in
\texttt{/tests/test\_holonomy.py}.

% ---------------------------------------------------------------------
\subsection*{3.9 Holonomy Drift Hallucinations (HDH)}

A latent trajectory exhibits HDH if there exists a sequence of loops
$\gamma_t$ such that
\[
\|H_{\gamma_t}\| > h_{\mathrm{thr}}
\quad \text{for infinitely many } t.
\]

Equivalently, HDH occurs when holonomy instability prevents the latent
trajectory from maintaining consistent parallel transport structure.

% =====================================================================
% End of 03-Holonomy.tex
% =====================================================================
